% ============================================================================
% sections/sec6_conclusion.tex
% ============================================================================
\section{Conclusion and Future Work}
\label{sec:conclusion}

\subsection{Conclusion}

A layered, empirically honest defense architecture against retrieval
poisoning was designed, implemented, and evaluated against real attack
constructions drawn from six cited papers. No single ring is sufficient
alone --- each has a real, measured limitation, stated in this report
rather than smoothed over --- but the combination of four rings plus a
persistent, cross-query trust store closes the measured gaps within this
simulation's scope, reaching $100.0\pm0.0\%$ accuracy and $0.0\pm0.0\%$
overall attack success across 1{,}600 evaluation queries
(\Cref{tab:main-comparison}).

Equally significant for a project of this kind is its debugging
discipline. Eight distinct, real, measured bugs were found across the
project's development --- from a self-referential filter threshold to a
consensus mechanism that silently did nothing at all
(\Cref{tab:debugging}) --- each fixed by direct measurement rather than
by adjusting a parameter until a target number appeared. That
verification trail is itself evidence of sound methodology and is
presented here alongside the final numbers rather than left out of the
report.

The most genuinely novel contributions of this work are, in order of how
directly they were earned through this project's own measurement rather
than inherited from the surveyed literature:

\begin{enumerate}
\item \textbf{Dual-Signal Risk Routing (\ringcode{claims/nli\_verifier.py})}:
  Combines geometric embedding cohesion with Natural Language Inference (NLI)
  premise-hypothesis contradiction scoring and date/numerical clash heuristics,
  discovered empirically during this project to overcome single-signal false
  consensus.
\item \textbf{Group-Wise Counterfactual Consensus \& Leave-Group-Out (LGO) Analysis (\ringcode{consensus/lgo\_analyzer.py})}:
  Generalizes RAGuard's leave-one-out to graph-theoretic community cliques,
  accurately isolating and quarantining colluding adversarial clusters.
\item \textbf{Chain-of-Verification (CoV) \& 5-State Calibrated Abstention (\ringcode{generation/})}:
  Proposition-level atomic claim cross-examination, cryptographic SHA-256 chunk
  citation auditing, and calibrated abstention (\texttt{ANSWER}, \texttt{PARTIAL\_ANSWER},
  \texttt{CONFLICTING\_EVIDENCE}, \texttt{INSUFFICIENT\_EVIDENCE}, \texttt{SECURITY\_BLOCK})
  to systematically prevent hallucination and sycophancy under adversarial pressure.
\item \textbf{Tenant-Scoped Dynamic Trust Ledger (\ringcode{trust/trust\_store.py})}:
  Closes the single-query information-theoretic ceiling by tracking persistent,
  decaying reputation scores across multi-turn sessions.
\item \textbf{Interactive Web Dashboard Studio (\texttt{dashboard/})}:
  A production-ready inspection workbench featuring real-time controlled poison
  injection sandboxing, live telemetry inspection drawers, side-by-side baseline
  comparisons, and local LLM backend integration (Ollama, OpenAI-compatible, and built-in).
\item \textbf{Unified Multi-Threat Statistical Harness (\texttt{benchmarks/})}:
  Evaluates 6 systems against 6 attack regimes with Student's-$t$ confidence
  intervals and an automated verification pipeline.
\end{enumerate}

\subsection{Future Work}
\label{sec:future-work}

\begin{enumerate}
\item \textbf{Dynamic Neural Vector Store Scaling}:
  Evaluate high-dimensional neural vector stores (e.g., Qdrant, Milvus, ChromaDB)
  with dense embeddings (e.g., BGE, E5, Contriever) under large-scale corpus
  scales ($10^5+$ documents).
\item \textbf{Adaptive Attacker Co-Evolution}:
  Test against adaptive, reinforcement-learning-guided attackers specifically
  optimized to evade Ring~0's repetition-ratio threshold and Ring~2's NLI
  contradiction heuristics simultaneously.
\item \textbf{Open-Domain Benchmark Integration}:
  Validate the complete defense pipeline on public multi-hop and adversarial RAG
  benchmarks (such as BEIR, HotpotQA, and Poisoned-RAG datasets) to further verify
  domain generalization across unstructured heterogeneous corpuses.
\item \textbf{Hardware Acceleration for Ring 3 Graph Partitions}:
  Explore GPU-accelerated sparse matrix graph clustering to scale Leave-Group-Out
  (LGO) consensus analysis to retrieval windows of $k > 50$ passages with sub-millisecond
  latencies.
\end{enumerate}
